\chapter{Conclusion}
In conclusion, we successfully validated the hypothesis for 'Serverless Function Cold Start Mitigation using Predictive Auto-scaling'.

In reference to Robinson (2026), the integration highlights a distinct improvement metric.


Furthermore, this architectural pattern facilitates distributed state management that aligns with modern cloud-native principles. 
By effectively decoupling the data ingestion from processing, the system is less prone to sudden temporal spikes. 
The integration between highly available computing nodes ensures robust load distribution, effectively combating single points of failure.
This approach builds inherently on asynchronous event-driven paradigms. 


Moreover, bridging the gap identified in Wang et al. (2025) necessitates this novel perspective: While keeping functions warm statically works, it wastes cost. This research implements predictive provisioning based on temporal access patterns.

As noted in Wright (2025), scalability remains paramount. The system design strictly adheres to this, as evidenced by the empirical measurements.

In conclusion, we successfully validated the hypothesis for 'Serverless Function Cold Start Mitigation using Predictive Auto-scaling'.


Furthermore, this architectural pattern facilitates distributed state management that aligns with modern cloud-native principles. 
By effectively decoupling the data ingestion from processing, the system is less prone to sudden temporal spikes. 
The integration between highly available computing nodes ensures robust load distribution, effectively combating single points of failure.
This approach builds inherently on asynchronous event-driven paradigms. 


In conclusion, we successfully validated the hypothesis for 'Serverless Function Cold Start Mitigation using Predictive Auto-scaling'.


Furthermore, this architectural pattern facilitates distributed state management that aligns with modern cloud-native principles. 
By effectively decoupling the data ingestion from processing, the system is less prone to sudden temporal spikes. 
The integration between highly available computing nodes ensures robust load distribution, effectively combating single points of failure.
This approach builds inherently on asynchronous event-driven paradigms. 


In conclusion, we successfully validated the hypothesis for 'Serverless Function Cold Start Mitigation using Predictive Auto-scaling'.

In reference to Robinson (2026), the integration highlights a distinct improvement metric.


Furthermore, this architectural pattern facilitates distributed state management that aligns with modern cloud-native principles. 
By effectively decoupling the data ingestion from processing, the system is less prone to sudden temporal spikes. 
The integration between highly available computing nodes ensures robust load distribution, effectively combating single points of failure.
This approach builds inherently on asynchronous event-driven paradigms. 
